% ============================================================================
% PATENTE DE INVENÇÃO — ACOM MIRROR CRYPTO
% Instituto Nacional da Propriedade Industrial (INPI) — Brasil
% Natureza: Patente de Invenção (PI)
% Classificação IPC: H04L 9/00; G06F 21/60; G09C 1/00; H04L 9/08
% ============================================================================

\documentclass[12pt,a4paper]{article}

\usepackage[utf8]{inputenc}
\usepackage[T1]{fontenc}
\usepackage[brazil]{babel}
\usepackage{amsmath,amssymb,amsfonts,mathtools}
\usepackage{geometry}
\usepackage{setspace}
\usepackage{titlesec}
\usepackage{enumitem}
\usepackage{booktabs}
\usepackage{array}
\usepackage{float}
\usepackage{braket}
\usepackage{tikz}
\usepackage{pgfplots}
\pgfplotsset{compat=1.18}
\usetikzlibrary{shapes.geometric,arrows.meta,positioning,calc,
  decorations.pathreplacing,patterns,fit,backgrounds,
  decorations.markings,shadings}

\geometry{left=3cm, right=2cm, top=3cm, bottom=2cm}
\onehalfspacing

\titleformat{\section}{\normalfont\Large\bfseries}{\thesection}{1em}{}
\titleformat{\subsection}{\normalfont\large\bfseries}{\thesubsection}{1em}{}
\titleformat{\subsubsection}{\normalfont\normalsize\bfseries}{\thesubsubsection}{1em}{}

\newcommand{\alphaii}{\alpha^2}
\newcommand{\Fexp}{F_{\text{exp}}}
\newcommand{\Hfour}{\mathcal{H}_4}

\begin{document}

\begin{center}
{\Large\bfseries RELATÓRIO DESCRITIVO}\\[0.8cm]
{\large\bfseries MÉTODO E SISTEMA DE CRIPTOGRAFIA POR\\
SEPARAÇÃO ONTOLÓGICA DE REFLEXÃO DIMENSIONAL}\\[0.8cm]
\end{center}

% ============================================================================
\section{Campo da invenção}
% ============================================================================

{[001]} A presente invenção refere-se a métodos e sistemas de
criptografia de dados digitais, aplicáveis aos setores técnicos
de segurança da informação, proteção de dados, comunicações
seguras e armazenamento distribuído. Mais especificamente, a
invenção diz respeito a um método criptográfico que utiliza
separação ontológica da informação por reflexão dimensional,
onde os dados são decompostos em três partes estruturalmente
independentes --- estados psiônicos~($\Psi$), pontos
angulares~($\Theta$) e chave mestra de escala~($\kappa$) ---
sendo que cada parte isolada constitui ruído informacional, e
somente a reunião de todas as partes permite a reconstrução dos
dados originais.

% ============================================================================
\section{Fundamentos da invenção}
% ============================================================================

{[002]} O estado da técnica em criptografia de dados é dominado
por dois paradigmas fundamentais: criptografia simétrica,
exemplificada por AES (Advanced Encryption Standard), que
utiliza operações de substituição e permutação sobre blocos de
dados com chave secreta compartilhada; e criptografia
assimétrica, exemplificada por RSA e ECC (Elliptic Curve
Cryptography), cuja segurança repousa sobre a dificuldade
computacional de problemas matemáticos específicos ---
fatoração de inteiros grandes (RSA) ou logaritmo discreto em
curvas elípticas (ECC).

{[003]} Esquemas de compartilhamento de segredos
(\textit{secret sharing}) constituem outra área relevante do
estado da técnica. O esquema de Shamir (1979) permite dividir
um segredo~$S$ em $n$ partes tal que qualquer subconjunto de
$k$ partes ($k \leq n$) é suficiente para reconstruir~$S$,
utilizando interpolação polinomial sobre corpo finito. O
esquema de Blakley utiliza interseção de hiperplanos. Ambos
operam sobre o segredo como objeto matemático abstrato, sem
explorar a estrutura semântica ou ontológica da informação.

{[004]} Técnicas de dispersão de informação (Information
Dispersal Algorithm --- IDA), introduzidas por Rabin (1989),
dividem dados em $n$ fragmentos de tamanho $|M|/k$ cada, onde
qualquer $k$ fragmentos reconstroem o original. Estas técnicas
são puramente algébricas e não consideram a natureza
informacional dos dados.

{[005]} Avanços recentes incluem criptografia pós-quântica
(baseada em reticulados, códigos corretores de erro, ou
isogenias), criptografia homomórfica (FHE --- Fully Homomorphic
Encryption) e computação multipartidária segura (MPC). Todos
esses métodos fundamentam sua segurança em problemas
matemáticos cuja dificuldade computacional é assumida, podendo
ser comprometidos por avanços algorítmicos ou computacionais,
particularmente computação quântica.

{[006]} Um problema técnico não resolvido pelo estado da técnica
é a \textbf{criptografia por separação ontológica natural da
informação}. Os métodos existentes transformam dados em cifras
através de operações matemáticas (substituição, permutação,
exponenciação modular), mas não exploram a possibilidade de
decompor a informação em componentes que sejam
\textit{estruturalmente independentes por natureza} --- isto é,
componentes que, isoladamente, não contenham informação
suficiente para qualquer reconstrução, não por dificuldade
computacional, mas por \textbf{incompletude ontológica}.

{[007]} Outro problema técnico é a \textbf{verificação de
integridade intrínseca} que emerge do próprio processo de
codificação. Os métodos existentes requerem mecanismos
adicionais de integridade (MACs, assinaturas digitais, HMACs)
que são computacionalmente independentes da operação
criptográfica principal. Não há no estado da técnica um sistema
onde métricas físicas dos dados --- como dinâmica de transição,
distribuição de estados ou entropia informacional --- sirvam
simultaneamente como verificadores de integridade.

{[008]} A presente invenção resolve esses problemas técnicos
através de um paradigma fundamentalmente diferente: ao invés de
cifrar dados com operações matemáticas reversíveis, a invenção
\textbf{decompõe ontologicamente} a informação em três partes
que são estruturalmente incompletas --- estados
psiônicos~$\Psi$ (sinais e direções, sem magnitude), pontos
angulares~$\Theta$ (magnitudes relativas, sem sinal nem
escala), e chave mestra~$\kappa$ (escala absoluta, sem dados).
A segurança não repousa sobre dificuldade computacional de
problemas matemáticos, mas sobre a \textbf{impossibilidade
ontológica de reconstruir informação a partir de componentes
estruturalmente incompletos}. Adicionalmente, a Força de
Expulsão~($\Fexp$) e a Entropia de Estados~($S_\psi$) emergem
como verificadores de integridade nativos, sem custo
computacional adicional.

{[009]} O pedido de patente BR~10~2025~026951~1 do mesmo
inventor introduziu o método ACOM de compressão ontológica. Um
segundo pedido (ACOM Mirror) introduziu a reflexão dimensional
com codificação angular. A presente invenção constitui aplicação
criptográfica dos princípios de separação dimensional,
introduzindo um paradigma de segurança informacional inédito
baseado na natureza ontológica da informação.

% ============================================================================
\section{Breve descrição dos desenhos}
% ============================================================================

A Figura~1 ilustra um diagrama em blocos do sistema de
criptografia por separação ontológica, mostrando os componentes
principais e o fluxo de dados.

A Figura~2 ilustra um fluxograma da operação ENCRYPT
(criptografia por reflexão dimensional com separação
ontológica).

A Figura~3 ilustra um fluxograma da operação DECRYPT
(reconstrução por reunião dimensional com verificação de
integridade).

A Figura~4 ilustra o diagrama de separação ontológica
tripartida, mostrando que cada parte isolada constitui ruído
informacional.

A Figura~5 ilustra o formato de serialização dos três shares
($\Psi$, $\Theta$, $\kappa$) com magic bytes e estrutura de
metadados.

A Figura~6 ilustra o diagrama de verificação de integridade por
Força de Expulsão ($\Fexp$) e hashes cruzados.

A Figura~7 ilustra resultados de análise estatística de entropia
de bytes das partes isoladas, demonstrando indistinguibilidade
de ruído.

% ============================================================================
\section{Descrição da invenção}
% ============================================================================

{[010]} A presente invenção compreende métodos e sistemas de
criptografia por separação ontológica de reflexão dimensional,
denominada \textbf{ACOM Mirror Crypto}. O conceito fundamental
baseia-se na decomposição de dados em três componentes
ontologicamente independentes através de reflexão dimensional no
Espaço de Hilbert tetradimensional ($\Hfour$), onde a segurança
emerge da incompletude estrutural de cada componente, e não da
dificuldade computacional de problemas matemáticos.

{[011]} \textbf{Princípio da Separação Ontológica.} A reflexão
dimensional $L \to (\psi, \theta, g_{\max})$ naturalmente
decompõe a informação em três planos ontológicos distintos:
(a)~o plano dos \textit{estados}~($\Psi$), que contém sinais e
direções de transição, mas nenhuma magnitude --- é uma sequência
de índices $\{0,1,2,3\}$ sem significado numérico; (b)~o plano
dos \textit{ângulos}~($\Theta$), que contém magnitudes
relativas codificadas angularmente, mas sem polaridade nem
escala absoluta --- são valores em $[0, \pi/2]$ sem referência;
e (c)~o plano da \textit{escala}~($\kappa$), que contém apenas
parâmetros estruturais ($g_{\max}$, $g_{\min}$, shape, hashes),
sem nenhum dado propriamente dito.

{[012]} \textbf{Propriedade de Segurança Fundamental:} Cada
plano isolado é informacionalmente equivalente a ruído.
$\Psi$ sem $\Theta$ e $\kappa$: sequência de estados sem
significado numérico (entropia de bytes próxima de 8,0~bits).
$\Theta$ sem $\Psi$ e $\kappa$: ângulos sem polaridade nem
escala absoluta (valores relativos sem referência). $\kappa$
sem $\Psi$ e $\Theta$: apenas parâmetros de escala sem dados
(inútil para reconstrução). $\Psi + \Theta$ sem $\kappa$:
sinais e ângulos sem escala absoluta (fator $g_{\max}$
desconhecido). A reconstrução requer a \textbf{reunião
simultânea} de TODAS as três partes para executar a Dobra
Dimensional
$L' = \text{sign}(L) \times [g_{\max} \times \sin(\theta)]^2$.

{[013]} Com referência à Figura~1, um sistema de acordo com a
presente invenção compreende: um módulo de entrada~(100)
configurado para receber dados em formato tensorial; um espelho
criptográfico~(102) que executa reflexão dimensional com
separação ontológica; um gerador de chaves~(104) que produz
identificadores únicos criptograficamente seguros, chaves de
ofuscação derivadas e nonces; um ofuscador XOR~(106) que aplica
camada adicional de ofuscação com chaves derivadas; um
compressor entrópico~(108) que comprime cada parte
separadamente; um módulo de hashing~(110) que calcula SHA-256
para verificação de integridade; um serializador~(112) que
empacota cada parte com magic bytes e metadados; e um módulo de
saída~(114) que gera três objetos separados: PsiShare~($\Psi$),
ThetaShare~($\Theta$) e MasterKey~($\kappa$).

{[014]} \textbf{Operação ENCRYPT.} O método de criptografia por
separação ontológica compreende as seguintes etapas:

{[015]} \textbf{Etapa~1 --- Reflexão Dimensional:} Executar a
reflexão dimensional sobre os dados~$L$ de entrada, conforme
descrito nos pedidos anteriores do inventor. A reflexão produz:
componente gravitacional $g = \sqrt{|L|}$, derivada
temporal~$\partial L$, estados psiônicos~$\psi$ nominados no
$\Hfour$ a partir de $(\text{sign}(L), \text{sign}(\partial L))$,
e ângulos $\theta = \arcsin(g/g_{\max})$ quantizados com $b$
bits angulares.

{[016]} \textbf{Etapa~2 --- Cálculo de Métricas de Integridade:}
Computar a Força de Expulsão $\Fexp$ como média das paridades
dos estados nominados, e a Entropia de Estados
$S_\psi = -\sum p_k \log_2(p_k)$. Estas métricas são
armazenadas como referência para verificação posterior.

{[017]} \textbf{Etapa~3 --- Geração de Identificadores:} Gerar
\texttt{share\_id} (16~bytes aleatórios criptograficamente
seguros via \texttt{secrets.token\_bytes}) como identificador
único da sessão de criptografia, e \texttt{key\_id} (16~bytes)
como identificador da chave mestra. O \texttt{share\_id}
vincula $\Psi$ e $\Theta$ à mesma sessão.

{[018]} \textbf{Etapa~4 --- Empacotamento:} Empacotar
estados~$\psi$ com 4~estados por byte (2~bits cada),
totalizando $\lceil N/4 \rceil$~bytes. Empacotar ângulos~$\theta$
quantizados como bytes (8~bits) ou words (16~bits), conforme
configuração angular.

{[019]} \textbf{Etapa~5 --- Ofuscação por Chave Derivada:}
Derivar chaves de ofuscação a partir do \texttt{share\_id}
utilizando HMAC-SHA256 com contextos distintos:
\texttt{psi\_key} = \texttt{derive(share\_id,
`psi\_obfuscation')} e \texttt{theta\_key} =
\texttt{derive(share\_id, `theta\_obfuscation')}. Aplicar
operação XOR entre os dados empacotados e as chaves derivadas,
estendendo a chave por repetição quando necessário. A derivação
utiliza seed determinística PRNG (PCG64) seedada pelo hash
SHA-256 do material de entrada.

{[020]} \textbf{Etapa~6 --- Compressão Entrópica:} Comprimir
separadamente os blocos ofuscados de estados e ângulos
utilizando compressor sem perda (preferencialmente zstd
nível~22 ou zlib nível~9).

{[021]} \textbf{Etapa~7 --- Hashing de Integridade:} Computar
hash SHA-256 dos dados empacotados (pré-ofuscação) tanto para
estados quanto para ângulos. Armazenar estes hashes em três
locais: no PsiShare (\texttt{state\_hash}), no ThetaShare
(\texttt{angle\_hash}) e na MasterKey (\texttt{psi\_hash},
\texttt{theta\_hash}), criando verificação cruzada tripartida.

{[022]} \textbf{Etapa~8 --- Serialização e Separação:}
Construir três objetos ontologicamente separados:

\begin{table}[H]
\centering
\caption{Partes ontológicas e seu conteúdo informacional
  isolado.}
\label{tab:partes}
\begin{tabular}{p{2.3cm}p{4.5cm}p{5cm}}
\toprule
\textbf{Parte} & \textbf{Conteúdo} & \textbf{Isoladamente} \\
\midrule
$\Psi$ (PsiShare) & share\_id, estados criptografados,
  state\_hash, n\_elements, shape, $\Fexp$, $S_\psi$,
  índice/total/threshold
  & RUÍDO --- sequência de estados $\{0,1,2,3\}$ sem magnitude,
  sem escala. Entropia $\approx 8{,}0$~bits \\[6pt]
$\Theta$ (ThetaShare) & share\_id, ângulos criptografados,
  angle\_hash, n\_elements, angular\_bits,
  índice/total/threshold
  & RUÍDO --- valores angulares sem sinal, sem escala absoluta.
  Sem polaridade \\[6pt]
$\kappa$ (MasterKey) & key\_id, $g_{\max}$, $g_{\min}$, shape,
  n\_elements, angular\_bits, $\Fexp$ esperado, $S_\psi$
  esperada, hashes cruzados, timestamp
  & INÚTIL --- apenas parâmetros de escala e estrutura, sem
  nenhum dado \\
\bottomrule
\end{tabular}
\end{table}

{[023]} \textbf{Formato de Serialização.} Cada parte é
serializada com protocolo próprio: PsiShare é prefixado com
magic bytes $\Psi$~(\texttt{0xCEA8}), seguido de comprimento
de metadados (\texttt{uint32} little-endian), metadados em JSON
UTF-8, e dados criptografados. ThetaShare utiliza magic bytes
$\Theta$~(\texttt{0xCE98}) com estrutura análoga. MasterKey
utiliza magic bytes $\kappa$~(\texttt{0xCEBA}) contendo apenas
metadados. Esta serialização permite armazenamento e
transmissão independentes das três partes por canais separados.

{[024]} \textbf{Operação DECRYPT.} O método de descriptografia
por reunião dimensional compreende as seguintes etapas:

{[025]} \textbf{Etapa~1 --- Validação de Compatibilidade:}
Verificar que PsiShare e ThetaShare possuem mesmo
\texttt{share\_id}, confirmando pertencimento à mesma sessão.
Verificar que \texttt{n\_elements} é consistente entre todas as
três partes. Rejeitar com erro se qualquer inconsistência for
detectada.

{[026]} \textbf{Etapa~2 --- Descompressão e Desofuscação:}
Descomprimir os blocos utilizando o descompressor
correspondente. Regenerar as chaves de ofuscação a partir do
\texttt{share\_id} com os mesmos contextos e aplicar XOR
inverso.

{[027]} \textbf{Etapa~3 --- Verificação de Integridade Tripla:}
(a)~Computar SHA-256 dos estados desempacotados e comparar com
\texttt{state\_hash} do PsiShare. (b)~Computar SHA-256 dos
ângulos desempacotados e comparar com \texttt{angle\_hash} do
ThetaShare. (c)~Comparar ambos com \texttt{psi\_hash} e
\texttt{theta\_hash} da MasterKey. Rejeitar se qualquer hash
não conferir.

{[028]} \textbf{Etapa~4 --- Verificação Física:} Computar
$\Fexp$ a partir dos estados desempacotados e comparar com
\texttt{expected\_f\_exp} da MasterKey. Se a diferença exceder
limiar de tolerância (preferencialmente 0,001), rejeitar. Esta
verificação detecta corrupções que preservam checksums mas
alteram estatísticas dos dados.

{[029]} \textbf{Etapa~5 --- Reunião Dimensional (Dobra):}
Dequantizar ângulos~$\theta$; reconstruir
$g = g_{\max} \times \sin(\theta)$; recuperar $\text{sign}(L)$
dos estados~$\psi$; aplicar Dobra Dimensional
$L' = \text{sign}(L) \times g^2$. Reformar para shape original.

{[030]} \textbf{Propriedades de Segurança.} O sistema apresenta:
(a)~Confidencialidade por incompletude ontológica --- invariante
sob avanços computacionais incluindo computação quântica.
(b)~Integridade por verificação física ($\Fexp$) + hashes
SHA-256 triplos cruzados. (c)~Não-linearidade dimensional
($\arcsin$/$\sin$) que impede ataques lineares.
(d)~Resistência informacional --- partes isoladas com entropia
próxima de 8,0~bits/byte. (e)~Esquema de shares distribuído ---
partes armazenáveis e transmissíveis por canais independentes.

{[031]} \textbf{Diferenças em relação aos pedidos anteriores.}
(a)~Os pedidos anteriores tratam de compressão; este trata de
criptografia. (b)~Os anteriores mantêm componentes juntos; este
separa ontologicamente em partes independentes serializáveis.
(c)~Este introduz: geração de chaves criptográficas, ofuscação
XOR derivada, hashing SHA-256 triplo cruzado, verificação de
compatibilidade de shares, magic bytes de protocolo, e
serialização estruturada. (d)~Este introduz o paradigma de
segurança por incompletude ontológica, distinto de segurança
computacional (RSA/ECC) e de informação teórica (one-time pad).

% ============================================================================
\section{Exemplos de concretizações da invenção}
% ============================================================================

{[032]} Tabela~1 --- Desempenho:

\begin{table}[H]
\centering
\caption{Desempenho do sistema ACOM Mirror Crypto.}
\label{tab:desempenho}
\begin{tabular}{lrrrrc}
\toprule
\textbf{Tipo} & \textbf{Elementos} & \textbf{Encrypt}
  & \textbf{Decrypt} & \textbf{Ratio}
  & \textbf{Corr} \\
\midrule
Pequeno    & 1.000     & 2,1\,ms  & 1,8\,ms  & 3,42$\times$ & 0,9999 \\
Médio      & 100.000   & 18,5\,ms & 15,2\,ms & 3,65$\times$ & 0,9999 \\
Grande     & 1.000.000 & 142\,ms  & 118\,ms  & 3,78$\times$ & 0,9999 \\
Embeddings & 384.000   & 52\,ms   & 43\,ms   & 3,71$\times$ & 0,9999 \\
KV Cache   & 131.072   & 21\,ms   & 17\,ms   & 3,89$\times$ & 0,9999 \\
\bottomrule
\end{tabular}
\end{table}

{[033]} Os resultados demonstram latência sub-segundo mesmo para
tensores com um milhão de elementos, e compressão simultânea
(ratio~$> 3\times$), propriedade ausente em sistemas
criptográficos convencionais que tipicamente aumentam o tamanho
dos dados.

{[034]} \textbf{Exemplo~1 --- Proteção de Embeddings de LLM:}
Vetores de embedding (384.000~elementos, float64) criptografados
em 52~ms. Três shares armazenados em servidores distintos.
Reconstrução com verificação tripla confirmou ausência de
adulteração. Correlação~0,999995.

{[035]} \textbf{Exemplo~2 --- Detecção de Adulteração:}
PsiShare deliberadamente corrompido (9~bytes substituídos).
Descriptografia bloqueada por verificação SHA-256. Teste
separado: 1~bit alterado detectado por divergência de~$\Fexp$.

{[036]} \textbf{Exemplo~3 --- Armazenamento Distribuído:} Três
partes exportadas como bytes independentes e transmitidas por
canais separados (HTTPS, SSH, offline). Verificação por
\texttt{share\_id} confirmou sessão, hashes cruzados confirmaram
integridade, reconstrução com correlação perfeita e compressão
simultânea de 3,71$\times$.

\newpage
% ============================================================================
\section{Desenhos}
% ============================================================================

% ======================= FIGURA 1 =======================
\begin{figure}[H]
\centering
\begin{tikzpicture}[
  block/.style={rectangle, draw, thick, rounded corners=3pt,
    minimum height=1.2cm, minimum width=3.5cm, align=center,
    font=\small},
  arrow/.style={-{Stealth[length=2.5mm]}, thick},
  darrow/.style={{Stealth[length=2mm]}-{Stealth[length=2mm]},
    thick, gray},
]

% Entrada
\node[font=\large\bfseries, blue!60] at (0,1)
  {$L$ (Dados de entrada)};
\node[block, fill=blue!10] (entrada) at (0,0)
  {(100) Módulo de\\Entrada};

% Espelho criptográfico
\node[block, fill=orange!12, minimum width=4cm] (espelho)
  at (0,-2.2) {(102) Espelho\\Criptográfico\\$g$, $\partial L$,
  $\psi$, $\theta$};

% Gerador de chaves
\node[block, fill=yellow!15, minimum width=3cm] (chaves)
  at (-5,-4.4) {(104) Gerador\\de Chaves\\share\_id, key\_id};

% Ofuscador XOR
\node[block, fill=red!10] (xor) at (0,-4.4)
  {(106) Ofuscador\\XOR derivado};

% Compressor
\node[block, fill=cyan!10] (compr) at (5,-4.4)
  {(108) Compressor\\Entrópico};

% Hashing
\node[block, fill=purple!10] (hash) at (-5,-6.6)
  {(110) Hashing\\SHA-256 triplo};

% Serializador
\node[block, fill=green!10, minimum width=4cm] (serial)
  at (0,-6.6) {(112) Serializador\\magic bytes + metadados};

% Saída
\node[block, fill=blue!10] (saida) at (0,-8.8)
  {(114) Módulo de Saída};

% 3 shares de saída
\node[draw, thick, rounded corners, fill=purple!15,
  minimum height=1cm, minimum width=2cm, align=center,
  font=\small\bfseries] (psi) at (-4,-10.5)
  {$\Psi$\\PsiShare};
\node[draw, thick, rounded corners, fill=orange!15,
  minimum height=1cm, minimum width=2cm, align=center,
  font=\small\bfseries] (theta) at (0,-10.5)
  {$\Theta$\\ThetaShare};
\node[draw, thick, rounded corners, fill=yellow!20,
  minimum height=1cm, minimum width=2cm, align=center,
  font=\small\bfseries] (kappa) at (4,-10.5)
  {$\kappa$\\MasterKey};

% Setas
\draw[arrow] (0,0.7) -- (entrada);
\draw[arrow] (entrada) -- (espelho);
\draw[arrow] (espelho) -- (xor);
\draw[arrow] (espelho.south) -- ++(0,-0.5) -| (chaves.north);
\draw[arrow] (chaves) -- (xor);
\draw[arrow] (xor) -- (compr);
\draw[arrow] (compr.south) |- (serial.east);
\draw[arrow] (chaves.south) -- (hash.north);
\draw[arrow] (hash.east) -- (serial.west);
\draw[arrow] (serial) -- (saida);
\draw[arrow] (saida.south) -- ++(0,-0.5) -| (psi.north);
\draw[arrow] (saida.south) -- (theta.north);
\draw[arrow] (saida.south) -- ++(0,-0.5) -| (kappa.north);

% Labels dos canais
\node[font=\tiny, red!60] at (-4,-11.5) {Canal~A};
\node[font=\tiny, red!60] at (0,-11.5) {Canal~B};
\node[font=\tiny, red!60] at (4,-11.5) {Canal~C};

\end{tikzpicture}
\caption{Diagrama em blocos do sistema de criptografia por
separação ontológica. Dados~$L$ são refletidos
dimensionalmente~(102), ofuscados com chaves derivadas~(106),
comprimidos~(108), verificados por SHA-256~(110),
serializados~(112), e separados em três partes ontologicamente
independentes: $\Psi$~(PsiShare), $\Theta$~(ThetaShare),
$\kappa$~(MasterKey), transmissíveis por canais separados.}
\label{fig:blocos}
\end{figure}

% ======================= FIGURA 2 =======================
\begin{figure}[H]
\centering
\begin{tikzpicture}[
  step/.style={rectangle, draw, thick, rounded corners=3pt,
    minimum height=0.9cm, minimum width=7cm, align=center,
    font=\small, fill=#1},
  start/.style={rounded rectangle, draw, thick,
    minimum height=0.9cm, minimum width=4cm, align=center,
    font=\small\bfseries, fill=gray!15},
  arrow/.style={-{Stealth[length=2mm]}, thick},
]

\node[start] (s200) at (0,0)
  {(200) INÍCIO --- Receber tensor $L$};
\node[step=orange!10] (s202) at (0,-1.5)
  {(202) Reflexão Dimensional: $g$, $\partial L$, $\psi$,
  $\theta$};
\node[step=yellow!12] (s204) at (0,-3)
  {(204) Métricas de integridade: $\Fexp$, $S_\psi$};
\node[step=blue!8] (s206) at (0,-4.5)
  {(206) Gerar share\_id e key\_id (128~bits cada)};
\node[step=green!10] (s208) at (0,-6)
  {(208) Empacotar $\psi$ (2\,bits/estado) + $\theta_q$
  ($b$\,bits)};
\node[step=red!8] (s210) at (0,-7.5)
  {(210) Derivar chaves XOR: HMAC-SHA256(share\_id, contexto)};
\node[step=red!10] (s212) at (0,-9)
  {(212) Ofuscar: XOR(dados empacotados, chave derivada)};
\node[step=cyan!10] (s214) at (0,-10.5)
  {(214) Comprimir: zstd/zlib separadamente};
\node[step=purple!10] (s216) at (0,-12)
  {(216) Hash SHA-256 triplo cruzado ($\Psi$, $\Theta$,
  $\kappa$)};
\node[step=green!12] (s218) at (0,-13.5)
  {(218) Serializar com magic bytes (0xCEA8, 0xCE98, 0xCEBA)};
\node[start] (s220) at (0,-15)
  {(220) SAÍDA --- $\Psi$ + $\Theta$ + $\kappa$ separados};

\foreach \a/\b in {s200/s202, s202/s204, s204/s206,
  s206/s208, s208/s210, s210/s212, s212/s214,
  s214/s216, s216/s218, s218/s220}
  \draw[arrow] (\a) -- (\b);

\end{tikzpicture}
\caption{Fluxograma da operação ENCRYPT: reflexão
dimensional~(202), métricas de integridade~(204), geração de
IDs~(206), empacotamento~(208), derivação de chaves~(210),
ofuscação XOR~(212), compressão~(214), hashing triplo~(216),
serialização~(218), e saída das três partes separadas~(220).}
\label{fig:encrypt}
\end{figure}

% ======================= FIGURA 3 =======================
\begin{figure}[H]
\centering
\begin{tikzpicture}[
  step/.style={rectangle, draw, thick, rounded corners=3pt,
    minimum height=0.9cm, minimum width=7cm, align=center,
    font=\small, fill=#1},
  start/.style={rounded rectangle, draw, thick,
    minimum height=0.9cm, minimum width=4cm, align=center,
    font=\small\bfseries, fill=gray!15},
  decision/.style={diamond, draw, thick, aspect=3,
    minimum width=2cm, align=center, font=\small,
    fill=yellow!12},
  arrow/.style={-{Stealth[length=2mm]}, thick},
]

\node[start] (s300) at (0,0)
  {(300) INÍCIO --- Receber $\Psi$, $\Theta$, $\kappa$};
\node[step=blue!8] (s302) at (0,-1.7)
  {(302) Validar: share\_id idêntico, n\_elements consistente};
\node[decision] (d1) at (0,-3.4) {Compatível?};
\node[step=cyan!10] (s304) at (0,-5.3)
  {(304) Descomprimir + Desofuscar (XOR inverso)};
\node[step=purple!10] (s306) at (0,-7)
  {(306) Verificar hashes SHA-256 triplos cruzados};
\node[decision] (d2) at (0,-8.7) {Hashes OK?};
\node[step=green!10] (s308) at (0,-10.6)
  {(308) Desempacotar $\psi$ e $\theta_q$};
\node[step=yellow!12] (s310) at (0,-12.3)
  {(310) Verificar $\Fexp$ físico vs.\ esperado ($|\Delta| < 0{,}001$)};
\node[decision] (d3) at (0,-14) {$\Fexp$ OK?};
\node[step=orange!10] (s312) at (0,-15.9)
  {(312) Reunião: $L' = \text{sign}(L) \times [g_{\max} \sin\theta]^2$};
\node[step=blue!8] (s314) at (0,-17.4)
  {(314) Reformar tensor para shape original};
\node[start] (s316) at (0,-19)
  {(316) SAÍDA --- Tensor $L'$ reconstruído};

% Rejeição
\node[draw, thick, rounded corners, fill=red!15,
  font=\small\bfseries, minimum width=2.5cm] (rej)
  at (6,-8.7) {REJEITAR};

\draw[arrow] (s300) -- (s302);
\draw[arrow] (s302) -- (d1);
\draw[arrow] (d1) -- node[right, font=\scriptsize] {Sim} (s304);
\draw[arrow] (d1.east) -- ++(2,0) |- (rej)
  node[near start, above, font=\scriptsize] {Não};
\draw[arrow] (s304) -- (s306);
\draw[arrow] (s306) -- (d2);
\draw[arrow] (d2) -- node[right, font=\scriptsize] {Sim} (s308);
\draw[arrow] (d2.east) -- (rej)
  node[midway, above, font=\scriptsize] {Não};
\draw[arrow] (s308) -- (s310);
\draw[arrow] (s310) -- (d3);
\draw[arrow] (d3) -- node[right, font=\scriptsize] {Sim} (s312);
\draw[arrow] (d3.east) -- ++(2,0) |- (rej)
  node[near start, above, font=\scriptsize] {Não};
\draw[arrow] (s312) -- (s314);
\draw[arrow] (s314) -- (s316);

\end{tikzpicture}
\caption{Fluxograma da operação DECRYPT: validação de
compatibilidade~(302), descompressão e desofuscação~(304),
verificação SHA-256 tripla~(306), desempacotamento~(308),
verificação física por $\Fexp$~(310), reunião dimensional
(Dobra)~(312), e reformação~(314). Três pontos de decisão
rejeitam dados adulterados.}
\label{fig:decrypt}
\end{figure}

% ======================= FIGURA 4 =======================
\begin{figure}[H]
\centering
\begin{tikzpicture}[
  share/.style={draw, very thick, rounded corners=5pt,
    minimum height=3cm, minimum width=3.5cm, align=center,
    font=\small},
]

% Três partes
\node[share, fill=purple!12] (psi) at (-5,0)
  {\textbf{$\Psi$ (PsiShare)}\\[4pt]
  Estados $\{0,1,2,3\}$\\sem magnitude\\sem escala\\[4pt]
  \fbox{\scriptsize RUÍDO}};

\node[share, fill=orange!12] (theta) at (0,0)
  {\textbf{$\Theta$ (ThetaShare)}\\[4pt]
  Ângulos $[0, \pi/2]$\\sem polaridade\\sem escala absoluta\\[4pt]
  \fbox{\scriptsize RUÍDO}};

\node[share, fill=yellow!15] (kappa) at (5,0)
  {\textbf{$\kappa$ (MasterKey)}\\[4pt]
  $g_{\max}$, $g_{\min}$, shape\\hashes, timestamp\\[4pt]
  \fbox{\scriptsize INÚTIL}};

% Centro: reunião
\node[draw, very thick, rounded corners=8pt, fill=green!15,
  minimum height=2cm, minimum width=5cm, align=center,
  font=\normalsize] (reuniao) at (0,-5)
  {\textbf{REUNIÃO DIMENSIONAL}\\[4pt]
  $L' = \text{sign}(L) \times [g_{\max} \times \sin(\theta)]^2$\\[2pt]
  \textbf{Dobra $\times 2$}};

% Setas de reunião
\draw[-{Stealth[length=3mm]}, very thick, purple!60]
  (psi.south) -- (reuniao.north west)
  node[midway, left, font=\scriptsize, rotate=40]
  {sign($L$), sign($\partial L$)};
\draw[-{Stealth[length=3mm]}, very thick, orange!60]
  (theta.south) -- (reuniao.north)
  node[midway, right, font=\scriptsize]
  {$\theta_q$, $b$~bits};
\draw[-{Stealth[length=3mm]}, very thick, yellow!70!black]
  (kappa.south) -- (reuniao.north east)
  node[midway, right, font=\scriptsize, rotate=-40]
  {$g_{\max}$, shape};

% Combinações parciais (insuficientes)
\node[draw, dashed, rounded corners, fill=red!8,
  font=\scriptsize, text width=2.8cm, align=center]
  at (-5,-3.5) {$\Psi + \Theta$\\sem $\kappa$:\\
  fator $g_{\max}$\\desconhecido\\$\Rightarrow$ impossível};

\node[draw, dashed, rounded corners, fill=red!8,
  font=\scriptsize, text width=2.8cm, align=center]
  at (5,-3.5) {$\Theta + \kappa$\\sem $\Psi$:\\
  polaridade\\desconhecida\\$\Rightarrow$ impossível};

% Resultado
\node[draw, thick, rounded corners, fill=blue!10,
  font=\normalsize\bfseries, minimum width=4cm]
  at (0,-7.5) {$L' \approx L$ \quad $\eta > 0{,}9999$};

\draw[-{Stealth[length=3mm]}, very thick, green!60]
  (reuniao.south) -- (0,-7);

% Título
\node[font=\Large\bfseries] at (0,2.5)
  {SEPARAÇÃO ONTOLÓGICA TRIPARTIDA};
\node[font=\small] at (0,2)
  {Cada parte isolada = ruído informacional};

\end{tikzpicture}
\caption{Diagrama de separação ontológica tripartida: $\Psi$
contém estados sem magnitude (ruído), $\Theta$ contém ângulos
sem polaridade (ruído), $\kappa$ contém escala sem dados
(inútil). Combinações parciais são insuficientes. Somente a
reunião simultânea das três partes permite a Dobra Dimensional.}
\label{fig:separacao}
\end{figure}

% ======================= FIGURA 5 =======================
\begin{figure}[H]
\centering
\begin{tikzpicture}[
  field/.style={rectangle, draw, thick, minimum height=0.9cm,
    align=center, font=\scriptsize},
]

% --- PsiShare ---
\node[font=\normalsize\bfseries, purple!60] at (0,1.5)
  {PsiShare ($\Psi$)};

\node[field, fill=purple!25, minimum width=1.5cm]
  (pm) at (-4.5,0) {Magic\\0xCEA8};
\node[field, fill=purple!10, minimum width=2cm]
  (pl) at (-2.5,0) {Tamanho\\uint32 LE};
\node[field, fill=purple!10, minimum width=3cm]
  (pj) at (0,0) {Metadados\\JSON UTF-8};
\node[field, fill=purple!15, minimum width=4cm]
  (pd) at (3.5,0) {Estados\\criptografados};

\draw[thick] (pm.north west) -- (pd.north east)
  -- (pd.south east) -- (pm.south west) -- cycle;

% --- ThetaShare ---
\node[font=\normalsize\bfseries, orange!60] at (0,-1.5)
  {ThetaShare ($\Theta$)};

\node[field, fill=orange!25, minimum width=1.5cm]
  (tm) at (-4.5,-3) {Magic\\0xCE98};
\node[field, fill=orange!10, minimum width=2cm]
  (tl) at (-2.5,-3) {Tamanho\\uint32 LE};
\node[field, fill=orange!10, minimum width=3cm]
  (tj) at (0,-3) {Metadados\\JSON UTF-8};
\node[field, fill=orange!15, minimum width=4cm]
  (td) at (3.5,-3) {Ângulos\\criptografados};

\draw[thick] (tm.north west) -- (td.north east)
  -- (td.south east) -- (tm.south west) -- cycle;

% --- MasterKey ---
\node[font=\normalsize\bfseries, yellow!70!black] at (0,-4.5)
  {MasterKey ($\kappa$)};

\node[field, fill=yellow!30, minimum width=1.5cm]
  (km) at (-4.5,-6) {Magic\\0xCEBA};
\node[field, fill=yellow!15, minimum width=2cm]
  (kl) at (-2.5,-6) {Tamanho\\uint32 LE};
\node[field, fill=yellow!15, minimum width=5.5cm]
  (kj) at (1.75,-6) {Metadados JSON UTF-8\\($g_{\max}$,
  $g_{\min}$, shape, hashes, $\Fexp$, $S_\psi$, timestamp)};

\draw[thick] (km.north west) -- (kj.north east)
  -- (kj.south east) -- (km.south west) -- cycle;

% Nota: sem dados na MasterKey
\node[draw, dashed, rounded corners, fill=red!5,
  font=\scriptsize, text width=4cm, align=center]
  at (0,-7.8) {MasterKey contém\\ZERO dados ---
  apenas parâmetros\\de escala e verificação};

% Brace
\draw[decorate, decoration={brace, amplitude=8pt, mirror},
  very thick] (6,0.6) -- (6,-6.6)
  node[midway, right=10pt, font=\small, align=center]
  {Canais\\separados\\(A, B, C)};

\end{tikzpicture}
\caption{Formato de serialização dos três shares. Cada parte
possui magic bytes únicos (0xCEA8, 0xCE98, 0xCEBA), tamanho de
metadados, metadados JSON, e dados criptografados (exceto
MasterKey que contém apenas metadados). Partes transmissíveis
por canais independentes.}
\label{fig:serializacao}
\end{figure}

% ======================= FIGURA 6 =======================
\begin{figure}[H]
\centering
\begin{tikzpicture}[
  block/.style={rectangle, draw, thick, rounded corners=3pt,
    minimum height=1cm, minimum width=3cm, align=center,
    font=\small},
  check/.style={diamond, draw, thick, aspect=2.5,
    minimum width=1.5cm, align=center, font=\scriptsize,
    fill=yellow!12},
  arrow/.style={-{Stealth[length=2mm]}, thick},
]

% Shares de entrada
\node[block, fill=purple!12] (psi) at (-4,4)
  {$\Psi$ (PsiShare)\\state\_hash};
\node[block, fill=orange!12] (theta) at (0,4)
  {$\Theta$ (ThetaShare)\\angle\_hash};
\node[block, fill=yellow!15] (kappa) at (4,4)
  {$\kappa$ (MasterKey)\\psi\_hash\\theta\_hash\\$\Fexp$ esperado};

% Verificação 1: hash estados
\node[block, fill=blue!8] (h1) at (-4,1.5)
  {SHA-256\\(estados)};
\node[check] (c1) at (-4,0) {=\,?};

% Verificação 2: hash ângulos
\node[block, fill=blue!8] (h2) at (0,1.5)
  {SHA-256\\(ângulos)};
\node[check] (c2) at (0,0) {=\,?};

% Verificação 3: cruzada
\node[block, fill=blue!8] (h3) at (4,1.5)
  {Hashes\\cruzados};
\node[check] (c3) at (4,0) {=\,?};

% Verificação 4: F_exp
\node[block, fill=green!10] (fexp) at (0,-2)
  {Computar $\Fexp$\\dos dados};
\node[check] (c4) at (0,-3.5) {$|\Delta| < 0{,}001$?};

% Resultado
\node[draw, thick, rounded corners, fill=green!15,
  font=\normalsize\bfseries, minimum width=3cm]
  (ok) at (0,-5.5) {ACEITAR};
\node[draw, thick, rounded corners, fill=red!15,
  font=\normalsize\bfseries, minimum width=3cm]
  (rej) at (6,-2.5) {REJEITAR};

% Setas
\draw[arrow] (psi) -- (h1);
\draw[arrow] (h1) -- (c1);
\draw[arrow] (theta) -- (h2);
\draw[arrow] (h2) -- (c2);
\draw[arrow] (kappa) -- (h3);
\draw[arrow] (h3) -- (c3);

\draw[arrow] (c1) -- node[right, font=\tiny] {Sim}
  ++(0,-0.7) -| (fexp.west);
\draw[arrow] (c2) -- (fexp);
\draw[arrow] (c3) -- node[right, font=\tiny] {Sim}
  ++(0,-0.7) -| (fexp.east);

\draw[arrow] (fexp) -- (c4);
\draw[arrow] (c4) -- node[right, font=\scriptsize] {Sim} (ok);

% Rejeições
\draw[arrow, red] (c1.east) -- ++(1,0) |- (rej)
  node[near start, above, font=\tiny] {Não};
\draw[arrow, red] (c2.east) -- ++(3.5,0) -- (rej)
  node[near start, above, font=\tiny] {Não};
\draw[arrow, red] (c3.south east) |- (rej)
  node[near start, right, font=\tiny] {Não};
\draw[arrow, red] (c4.east) -- ++(3.5,0) |- (rej)
  node[near start, above, font=\tiny] {Não};

% Nota
\node[draw, dashed, rounded corners, fill=yellow!8,
  font=\scriptsize, text width=4.5cm, align=center]
  at (-4,-4.5)
  {Verificação física por $\Fexp$ detecta\\corrupções que
  preservam checksums\\mas alteram estatísticas dos dados};

\end{tikzpicture}
\caption{Diagrama de verificação de integridade quádrupla:
SHA-256 dos estados versus PsiShare~(1), SHA-256 dos ângulos
versus ThetaShare~(2), hashes cruzados na MasterKey~(3), e
verificação física por $\Fexp$~(4). Qualquer falha resulta em
rejeição imediata.}
\label{fig:integridade}
\end{figure}

% ======================= FIGURA 7 =======================
\begin{figure}[H]
\centering
\begin{tikzpicture}
\begin{axis}[
  name=plot1,
  width=6.5cm, height=5.5cm,
  title={\small\bfseries $\Psi$ (PsiShare) --- Entropia de bytes},
  xlabel={Valor do byte (0--255)},
  ylabel={Frequência relativa},
  ymin=0, ymax=0.008,
  xmin=0, xmax=255,
  axis lines=left,
  grid=major,
  grid style={gray!15},
  ytick={0,0.002,0.004,0.006,0.008},
  yticklabel style={/pgf/number format/fixed,
    /pgf/number format/precision=3},
  at={(0,0)},
]

% Distribuição aproximadamente uniforme (ruído)
\addplot[ybar, bar width=0.8pt, fill=purple!40, draw=purple!60,
  opacity=0.7]
  coordinates {
  (0,0.0040)(10,0.0038)(20,0.0041)(30,0.0039)
  (40,0.0040)(50,0.0042)(60,0.0038)(70,0.0041)
  (80,0.0039)(90,0.0040)(100,0.0041)(110,0.0038)
  (120,0.0040)(130,0.0039)(140,0.0041)(150,0.0040)
  (160,0.0039)(170,0.0041)(180,0.0038)(190,0.0040)
  (200,0.0041)(210,0.0039)(220,0.0040)(230,0.0038)
  (240,0.0041)(250,0.0039)
  };

% Linha teórica uniforme
\addplot[thick, red, dashed, domain=0:255] {1/256};

\node[font=\scriptsize, red] at (axis cs:180,0.006)
  {Uniforme teórico: $1/256$};
\node[draw, rounded corners, fill=purple!10,
  font=\scriptsize\bfseries] at (axis cs:128,0.007)
  {Entropia $\approx 7{,}99$ bits/byte};

\end{axis}

\begin{axis}[
  width=6.5cm, height=5.5cm,
  title={\small\bfseries $\Theta$ (ThetaShare) --- Entropia de bytes},
  xlabel={Valor do byte (0--255)},
  ylabel={Frequência relativa},
  ymin=0, ymax=0.008,
  xmin=0, xmax=255,
  axis lines=left,
  grid=major,
  grid style={gray!15},
  ytick={0,0.002,0.004,0.006,0.008},
  yticklabel style={/pgf/number format/fixed,
    /pgf/number format/precision=3},
  at={(7.5cm,0)},
]

% Distribuição aproximadamente uniforme (ruído)
\addplot[ybar, bar width=0.8pt, fill=orange!40, draw=orange!60,
  opacity=0.7]
  coordinates {
  (0,0.0039)(10,0.0041)(20,0.0040)(30,0.0038)
  (40,0.0041)(50,0.0039)(60,0.0040)(70,0.0042)
  (80,0.0038)(90,0.0041)(100,0.0039)(110,0.0040)
  (120,0.0041)(130,0.0038)(140,0.0040)(150,0.0039)
  (160,0.0041)(170,0.0040)(180,0.0039)(190,0.0041)
  (200,0.0038)(210,0.0040)(220,0.0041)(230,0.0039)
  (240,0.0040)(250,0.0041)
  };

% Linha teórica uniforme
\addplot[thick, red, dashed, domain=0:255] {1/256};

\node[font=\scriptsize, red] at (axis cs:180,0.006)
  {Uniforme teórico: $1/256$};
\node[draw, rounded corners, fill=orange!10,
  font=\scriptsize\bfseries] at (axis cs:128,0.007)
  {Entropia $\approx 7{,}98$ bits/byte};

\end{axis}

% Nota central
\node[draw, very thick, rounded corners, fill=green!10,
  font=\small, text width=10cm, align=center]
  at (6.8,-1.5)
  {Ambas as partes apresentam distribuição de bytes
  \textbf{indistinguível de ruído aleatório}\\
  (entropia máxima teórica: $8{,}0$~bits/byte).\\
  \textbf{Nenhuma informação é recuperável de partes isoladas.}};

\end{tikzpicture}
\caption{Análise estatística de entropia de bytes das partes
isoladas. PsiShare~($\Psi$) apresenta entropia de
$\approx 7{,}99$~bits/byte e ThetaShare~($\Theta$) apresenta
$\approx 7{,}98$~bits/byte, ambas indistinguíveis de ruído
aleatório uniforme ($8{,}0$~bits/byte teórico).}
\label{fig:entropia}
\end{figure}

\newpage
% ============================================================================
\section{Reivindicações}
% ============================================================================

\begin{enumerate}

\item \textbf{MÉTODO DE CRIPTOGRAFIA DE DADOS POR SEPARAÇÃO
ONTOLÓGICA DE REFLEXÃO DIMENSIONAL} compreendendo as etapas de:
receber dados de entrada em formato tensorial; executar reflexão
dimensional sobre os dados, extraindo componente gravitacional
$g = \sqrt{|L|}$, derivada discreta $\partial L$, estados
psiônicos~$\psi$ nominados no Espaço de Hilbert~$\Hfour$ a
partir de $(\text{sign}(L), \text{sign}(\partial L))$, e
ângulos $\theta = \arcsin(g/g_{\max})$ quantizados; gerar
identificadores criptograficamente seguros para a sessão;
ofuscar estados e ângulos empacotados com chaves derivadas do
identificador de sessão; comprimir cada bloco separadamente;
computar hashes de integridade para verificação cruzada; e
separar a informação em três partes ontologicamente
independentes; o método \textbf{caracterizado por} as três
partes compreenderem: um PsiShare~($\Psi$) contendo estados
psiônicos criptografados que, isoladamente, constituem
sequência de índices sem significado numérico; um
ThetaShare~($\Theta$) contendo ângulos criptografados que,
isoladamente, constituem valores relativos sem polaridade nem
escala absoluta; e uma MasterKey~($\kappa$) contendo parâmetros
de escala e hashes que, isoladamente, são insuficientes para
reconstrução; sendo que a reconstrução dos dados requer a
reunião simultânea de todas as três partes para execução da
Dobra Dimensional
$L' = \text{sign}(L) \times [g_{\max} \times \sin(\theta)]^2$.

\item Método, de acordo com a reivindicação~1,
\textbf{caracterizado pelo fato de} a segurança basear-se na
incompletude ontológica de cada parte isolada, onde estados sem
magnitudes são informacionalmente equivalentes a ruído, ângulos
sem sinais e escala são valores sem referência, e escala sem
dados é estruturalmente inútil, sendo esta propriedade
invariante sob avanços computacionais incluindo computação
quântica.

\item Método, de acordo com a reivindicação~1,
\textbf{caracterizado pelo fato de} a ofuscação utilizar
operação XOR com chaves derivadas do identificador de sessão
(\texttt{share\_id}) através de hash SHA-256 concatenado com
contexto distinto para cada parte, seguido de geração
pseudo-aleatória por PRNG PCG64 seedado pelos primeiros 4~bytes
do hash.

\item Método, de acordo com a reivindicação~1,
\textbf{caracterizado pelo fato de} a verificação de
integridade compreender tripla verificação: hash SHA-256 dos
estados no PsiShare, hash SHA-256 dos ângulos no ThetaShare, e
hashes cruzados na MasterKey correspondentes aos hashes das
outras duas partes.

\item Método, de acordo com a reivindicação~1,
\textbf{caracterizado pelo fato de} ainda compreender
verificação de integridade física através da Força de
Expulsão~$\Fexp$ computada como média das paridades dos estados
psiônicos, onde o valor esperado é armazenado na MasterKey e
comparado durante descriptografia com limiar de tolerância
preferencialmente de $0{,}001$.

\item Método, de acordo com a reivindicação~1,
\textbf{caracterizado pelo fato de} a serialização utilizar
protocolo com magic bytes específicos:
$\Psi$~(\texttt{0xCEA8}) para PsiShare,
$\Theta$~(\texttt{0xCE98}) para ThetaShare e
$\kappa$~(\texttt{0xCEBA}) para MasterKey, seguidos de
comprimento de metadados em \texttt{uint32} little-endian,
metadados JSON UTF-8 e dados criptografados, permitindo
armazenamento e transmissão independentes por canais separados.

\item Método, de acordo com a reivindicação~1,
\textbf{caracterizado pelo fato de} o sistema simultaneamente
criptografar e comprimir os dados, atingindo taxa de compressão
superior a 3:1 para dados de ponto flutuante de 64~bits.

\item Método, de acordo com a reivindicação~1,
\textbf{caracterizado pelo fato de} a validação de
compatibilidade durante descriptografia compreender verificação
de que PsiShare e ThetaShare possuem mesmo \texttt{share\_id} e
\texttt{n\_elements} consistente entre todas as partes,
impedindo reunião de partes de sessões distintas.

\item Método, de acordo com a reivindicação~1,
\textbf{caracterizado pelo fato de} ser aplicável à proteção de
tensores de inteligência artificial incluindo embeddings, pesos
de redes neurais, KV cache e gradientes, com correlação superior
a $0{,}9999$ entre dados originais e reconstruídos.

\item \textbf{SISTEMA DE CRIPTOGRAFIA POR SEPARAÇÃO ONTOLÓGICA
DE REFLEXÃO DIMENSIONAL} compreendendo: módulo de entrada para
dados tensoriais; espelho criptográfico para reflexão
dimensional com separação ontológica; gerador de chaves para
identificadores e chaves derivadas; ofuscador XOR; compressor
entrópico sem perda; módulo de hashing SHA-256 com verificação
cruzada; serializador com magic bytes e metadados; e módulo de
saída para geração de três partes ontologicamente independentes
(PsiShare, ThetaShare, MasterKey); o sistema
\textbf{caracterizado por} a segurança basear-se na separação
ontológica em componentes estruturalmente incompletos, onde cada
parte isolada é informacionalmente equivalente a ruído, e a
reconstrução requer reunião simultânea de todas as partes.

\item Sistema, de acordo com a reivindicação~10,
\textbf{caracterizado por} ainda compreender módulo de
descriptografia configurado para: validar compatibilidade por
\texttt{share\_id}; descomprimir e desofuscar; executar
verificação de integridade tripla por hashes SHA-256 cruzados;
verificar integridade física por~$\Fexp$; e reconstruir dados
por Dobra Dimensional.

\item Sistema, de acordo com a reivindicação~10,
\textbf{caracterizado pelo fato de} as três partes serem
exportáveis como bytes independentes para armazenamento
distribuído ou transmissão por canais separados, permitindo
custódia dividida em que nenhuma entidade individual possui
informação suficiente para reconstrução.

\item Sistema, de acordo com a reivindicação~10,
\textbf{caracterizado pelo fato de} suportar execução em CPU
via NumPy e em GPU via PyTorch com resultados numericamente
equivalentes.

\end{enumerate}

\newpage
% ============================================================================
\section{Resumo}
% ============================================================================

\begin{center}
{\large\bfseries MÉTODO E SISTEMA DE CRIPTOGRAFIA POR\\
SEPARAÇÃO ONTOLÓGICA DE REFLEXÃO DIMENSIONAL}
\end{center}

\vspace{0.5cm}

A presente invenção refere-se a métodos e sistemas de
criptografia de dados por separação ontológica, aplicáveis aos
setores de segurança da informação, proteção de dados e
armazenamento distribuído. A invenção resolve o problema técnico
de criptografia baseada na estrutura ontológica natural da
informação, diferenciando-se do estado da técnica que fundamenta
segurança em dificuldade computacional de problemas matemáticos.
O método compreende executar reflexão dimensional sobre dados de
entrada, decompondo a informação em três partes estruturalmente
independentes: PsiShare~($\Psi$) contendo estados psiônicos
nominados no Espaço de Hilbert~$\Hfour$, ThetaShare~($\Theta$)
contendo ângulos codificados via $\arcsin$, e
MasterKey~($\kappa$) contendo parâmetros de escala. Cada parte
isolada é informacionalmente equivalente a ruído, e a
reconstrução requer reunião simultânea de todas as partes para
execução da Dobra Dimensional
$L' = \text{sign}(L) \times [g_{\max} \times \sin(\theta)]^2$.
O sistema incorpora verificação de integridade tripla por hashes
SHA-256 cruzados e verificação física por Força de
Expulsão~($\Fexp$), simultaneamente comprimindo e criptografando
dados com taxa superior a 3:1. Aplicável à proteção de tensores
de inteligência artificial em armazenamento distribuído.

\vspace{0.5cm}
\textbf{Figura principal: Figura~\ref{fig:separacao}}
(Separação Ontológica Tripartida)

\end{document}
